\section{Introduction}
A transformer attends differently to different prompts. Intuition
suggests that a maximally random prompt forces every token to attend
broadly, while a highly compressible prompt admits attention that
concentrates on a small number of pattern-matching keys. We make this
intuition quantitative.

\paragraph{Question.} Is there a measurable, monotone relationship
between the algorithmic complexity of the prompt and the density of
attention used to process it?

Kolmogorov complexity $K(x)$ is uncomputable, but it admits an
upper-bound proxy via any reasonable lossless compressor:
$K(x) \le |\text{gzip}(x)| + O(1)$~\cite{li2008introduction}. The
gzip-ratio
$\rho(x) := |\text{gzip}(x)|/|x|$ has been used as a complexity proxy
in many recent works~\cite{jiang2023low,delete2024gzip}.

\paragraph{Findings.} Across seven prompt families and GPT-2's twelve
layers, mean attention entropy is well-predicted by $\rho(x)$ in the
mid-network: at layer $\gptTwoMaxRsqLayer$ a univariate linear
regression achieves $R^2=\gptTwoMaxRsq$. The relationship is weaker
in the embedding-near and deepest layers, suggesting a U-shaped
dependence of attention adaptability on depth.

\paragraph{Contributions.} (i) A reproducible pipeline mapping
compressibility to per-layer attention-density metrics (entropy, Gini,
sparsity); (ii) a per-layer regression analysis isolating the depth
range in which the relationship holds; (iii) an open dataset of
controlled prompt families spanning the algorithmic-complexity spectrum.
